% =====================================================================
% スペクトル指数ランダマイゼーションによる色付き観測ノイズ
% Sim-to-Realドメインランダマイゼーション (Overleaf対応版)
%
% 必要なパッケージ:
%   \usepackage{amsmath}
%   \usepackage{amssymb}
%   \usepackage{bm}
%   \usepackage{booktabs}
% =====================================================================

\subsection{スペクトル指数ランダマイゼーションによる色付き観測ノイズ}
\label{sec:colored_obs_noise}

\subsubsection{問題設定}

ドメインランダマイゼーションでは，Sim-to-Realギャップを埋めるために
シミュレーション上の観測にノイズを付加する手法が広く用いられている．
従来は，各ステップで独立同分布（i.i.d.）の白色ノイズ
$\boldsymbol{\epsilon}_t \sim \mathcal{N}(\boldsymbol{0}, \sigma^2 \boldsymbol{I})$
を付加することが一般的であった．
しかし，実機のセンサ誤差には時間的相関が存在する．
エンコーダのバックラッシュ，サーモドリフト，通信ジッタは
それぞれ固有の周波数スペクトルを持つノイズを発生させる．
この課題に対し，時間相関を持つピンクノイズ（$1/f$ノイズ，スペクトル指数$\beta=1$）を
用いる手法が提案されている~\cite{eberhard2023pink}．

しかし，既存手法ではスペクトル指数が$\beta = 1$に固定されている点に限界がある．
実際のロボットでは，異なるノイズ源が異なるスペクトル構造を示す：
エンコーダの量子化誤差は$\beta \approx 0.5$，
サーボのバックラッシュは$\beta \approx 1.0$--$1.5$，
サーモドリフトは$\beta \approx 2.0$に近い特性を持つ．
単一の固定$\beta$で学習されたポリシーは，
その特定のノイズ構造に過適合するリスクがある．

\subsubsection{提案手法：エピソードごとの$\beta$ランダマイゼーション}

本研究では，スペクトル指数$\beta$をドメインランダマイゼーションのパラメータとして
ランダマイズすることを提案する．
各エピソードのリセット時に，各並列シミュレーション環境が独立に
\begin{equation}
  \beta_i \sim \mathcal{U}(\beta_{\min},\, \beta_{\max})
  \label{eq:beta_sample}
\end{equation}
をサンプリングする．
ここで$\beta_{\min}$と$\beta_{\max}$は，学習時に経験するノイズスペクトルの範囲を定義する．
サンプリングされた$\beta_i$に基づき，パワースペクトル密度
$S(f) \propto 1/f^{\beta_i}$に従う色付きノイズ系列を
Timmer--Koenig FFTアルゴリズム~\cite{timmer1995}で生成する．
従来のピンクノイズ生成との差異は，
周波数領域における振幅スケーリングがサンプリングされた指数を用いる点にある：
\begin{equation}
  s_k = f_k^{-\beta_i / 2}
  \label{eq:scaling}
\end{equation}
これは標準的なピンクノイズで用いられる固定値$s_k = f_k^{-1/2}$を一般化したものである．

生成された単位分散のノイズは，センサ由来の観測（関節角度・関節角速度）にのみ選択的に注入し，
指令値や前ステップアクションにはノイズを付加しない：
\begin{equation}
  \tilde{\boldsymbol{o}}_t^{(\mathrm{sens})} = \boldsymbol{o}_t^{(\mathrm{sens})} + \sigma_{\mathrm{obs}} \, \boldsymbol{\eta}_t^{(\beta_i)}
  \label{eq:corruption}
\end{equation}
ここで$\boldsymbol{\eta}_t^{(\beta_i)}$は指数$\beta_i$で生成された色付きノイズサンプルである．

\subsubsection{比較条件}

$\beta$ランダマイゼーションの効果を検証するため，
既存のベースラインに加えて以下の比較条件を設定した．

\begin{table}[htbp]
  \centering
  \caption{ノイズ注入条件の比較}
  \label{tab:colored_scenarios}
  \begin{tabular}{lccc}
    \toprule
    条件名 & アクションノイズ & 観測ノイズ & $\beta$ \\
    \midrule
    No-Noise（ベースライン） & -- & -- & -- \\
    Obs-Noise（ピンク） & -- & ピンク & $1$（固定） \\
    \textbf{Obs-Colored-Noise（提案手法）} & -- & \textbf{色付き} & $\boldsymbol{\mathcal{U}(0.5, 1.5)}$ \\
    Both-Colored-Noise & ピンク & 色付き & アクション: $1$，観測: $\mathcal{U}(0.5, 1.5)$ \\
    \bottomrule
  \end{tabular}
\end{table}

\subsubsection{ハイパーパラメータ}

\begin{table}[htbp]
  \centering
  \caption{色付き観測ノイズのハイパーパラメータ}
  \label{tab:colored_hyperparams}
  \begin{tabular}{lll}
    \toprule
    パラメータ & 記号 & 値 \\
    \midrule
    スペクトル指数の範囲 & $[\beta_{\min},\, \beta_{\max}]$ & $[0.5,\, 1.5]$ \\
    観測ノイズスケール & $\sigma_{\mathrm{obs}}$ & $0.02$ \\
    ノイズバッファ長 & $T$ & $1000$ ステップ \\
    \bottomrule
  \end{tabular}
\end{table}

$\beta \in [0.5, 1.5]$の範囲は，
弱い相関を持つほぼ白色のノイズ（$\beta = 0.5$）から
強い時間相関を持つノイズ（$\beta = 1.5$，ブラウンノイズに近い）までをカバーし，
実機センサの多様な特性を網羅する設定となっている．

% =====================================================================
% 参考文献（.bib ファイルに追加してください）
% =====================================================================
% @inproceedings{eberhard2023pink,
%   title     = {Pink Noise Is All You Need: Colored Noise Exploration in Deep Reinforcement Learning},
%   author    = {Eberhard, Onno and Hollenstein, Jakob and Pinneri, Cristina and Martius, Georg},
%   booktitle = {Proceedings of the Eleventh International Conference on Learning Representations (ICLR 2023)},
%   year      = {2023},
%   url       = {https://openreview.net/forum?id=hQ9V5QN27eS},
% }
% @article{timmer1995,
%   title   = {On generating power law noise},
%   author  = {Timmer, J. and Koenig, M.},
%   journal = {Astronomy and Astrophysics},
%   volume  = {300},
%   pages   = {707--710},
%   year    = {1995},
% }
